\documentclass[11pt]{article}
% The preamble every project's paper shares: % The preamble every project's paper shares: % The preamble every project's paper shares: \input{../../preamble} from
% projects/<project>/paper/. Boards come from chessboard and xskak, plots from
% pgfplots; every number a paper prints comes from its generated.tex.
\usepackage[LSBC4,T1]{fontenc}
\usepackage[utf8]{inputenc}
\usepackage{lmodern}
\usepackage[margin=1in]{geometry}
\usepackage{microtype}
\usepackage{amsmath}
\usepackage{booktabs}
\usepackage{tabularx}
\usepackage{array}
\usepackage{xcolor}
\usepackage[obeyspaces,spaces]{url}
\usepackage{xskak}
\usepackage{pgfplots}
\pgfplotsset{compat=1.18}
\usepackage{caption}
\usepackage{subcaption}
\usepackage{enumitem}
\usepackage[section]{placeins}
\usepackage[hidelinks]{hyperref}

\setchessboard{boardfontencoding=LSBC4, boardfontfamily=skaknew, blackfieldcolor=black!22,
  whitefieldcolor=white, setfontcolors, showmover=false, label=true, labelfontsize=6pt,
  boardfontsize=14pt}

% A move or line in UCI, and a position in FEN, both breakable at a slash.
\newcommand{\uci}[1]{\texttt{#1}}
\newcommand{\fen}[1]{{\footnotesize\path{#1}}}
\newcolumntype{L}{>{\raggedright\arraybackslash}X}

\setlength{\parskip}{0.4em}
\setlength{\parindent}{0pt}
 from
% projects/<project>/paper/. Boards come from chessboard and xskak, plots from
% pgfplots; every number a paper prints comes from its generated.tex.
\usepackage[LSBC4,T1]{fontenc}
\usepackage[utf8]{inputenc}
\usepackage{lmodern}
\usepackage[margin=1in]{geometry}
\usepackage{microtype}
\usepackage{amsmath}
\usepackage{booktabs}
\usepackage{tabularx}
\usepackage{array}
\usepackage{xcolor}
\usepackage[obeyspaces,spaces]{url}
\usepackage{xskak}
\usepackage{pgfplots}
\pgfplotsset{compat=1.18}
\usepackage{caption}
\usepackage{subcaption}
\usepackage{enumitem}
\usepackage[section]{placeins}
\usepackage[hidelinks]{hyperref}

\setchessboard{boardfontencoding=LSBC4, boardfontfamily=skaknew, blackfieldcolor=black!22,
  whitefieldcolor=white, setfontcolors, showmover=false, label=true, labelfontsize=6pt,
  boardfontsize=14pt}

% A move or line in UCI, and a position in FEN, both breakable at a slash.
\newcommand{\uci}[1]{\texttt{#1}}
\newcommand{\fen}[1]{{\footnotesize\path{#1}}}
\newcolumntype{L}{>{\raggedright\arraybackslash}X}

\setlength{\parskip}{0.4em}
\setlength{\parindent}{0pt}
 from
% projects/<project>/paper/. Boards come from chessboard and xskak, plots from
% pgfplots; every number a paper prints comes from its generated.tex.
\usepackage[LSBC4,T1]{fontenc}
\usepackage[utf8]{inputenc}
\usepackage{lmodern}
\usepackage[margin=1in]{geometry}
\usepackage{microtype}
\usepackage{amsmath}
\usepackage{booktabs}
\usepackage{tabularx}
\usepackage{array}
\usepackage{xcolor}
\usepackage[obeyspaces,spaces]{url}
\usepackage{xskak}
\usepackage{pgfplots}
\pgfplotsset{compat=1.18}
\usepackage{caption}
\usepackage{subcaption}
\usepackage{enumitem}
\usepackage[section]{placeins}
\usepackage[hidelinks]{hyperref}

\setchessboard{boardfontencoding=LSBC4, boardfontfamily=skaknew, blackfieldcolor=black!22,
  whitefieldcolor=white, setfontcolors, showmover=false, label=true, labelfontsize=6pt,
  boardfontsize=14pt}

% A move or line in UCI, and a position in FEN, both breakable at a slash.
\newcommand{\uci}[1]{\texttt{#1}}
\newcommand{\fen}[1]{{\footnotesize\path{#1}}}
\newcolumntype{L}{>{\raggedright\arraybackslash}X}

\setlength{\parskip}{0.4em}
\setlength{\parindent}{0pt}

% Written by projects/maia-punish/paper/generate.mjs; do not edit by hand.

\newcommand{\StepAt}{10}
\newcommand{\ErrorAt}{5}
\newcommand{\EAt}{2}
\newcommand{\BinMin}{15}
\newcommand{\BinGap}{0.1}
\newcommand{\Draws}{1000}
\newcommand{\Events}{689}
\newcommand{\ExactEvents}{517}
\newcommand{\Games}{60}
\newcommand{\BandMin}{1000}
\newcommand{\BandMax}{1400}
\newcommand{\ResultTable}{%
\begin{tabular}{llrrrrlll}
\toprule
replies & errors of & events & punished & $q$ & base & size & $q$ $-$ base & $q$ $-$ size \\
\midrule
verified & player & 243 (57) & 0.436 & 0.162 & 0.249 & 0.236 & $-$0.087 [$-$0.118, $-$0.056] & $-$0.074 [$-$0.106, $-$0.046] \\
verified & opponents & 274 (59) & 0.427 & 0.145 & 0.247 & 0.231 & $-$0.102 [$-$0.125, $-$0.079] & $-$0.087 [$-$0.113, $-$0.060] \\
all & player & 334 (57) & 0.419 & 0.162 & 0.245 & 0.241 & $-$0.083 [$-$0.105, $-$0.060] & $-$0.079 [$-$0.103, $-$0.055] \\
all & opponents & 355 (60) & 0.425 & 0.151 & 0.246 & 0.238 & $-$0.095 [$-$0.116, $-$0.074] & $-$0.087 [$-$0.110, $-$0.066] \\
\bottomrule
\end{tabular}
}
\newcommand{\PlayerEvents}{243}
\newcommand{\PlayerGames}{57}
\newcommand{\PlayerPunished}{0.44}
\newcommand{\PlayerBrier}{0.162}
\newcommand{\PlayerBrierBase}{0.249}
\newcommand{\PlayerBrierSize}{0.236}
\newcommand{\PlayerBigErrors}{140}
\newcommand{\PlayerUnpunished}{45}
\newcommand{\PlayerQPunished}{0.62}
\newcommand{\PlayerQUnpunished}{0.28}
\newcommand{\PlayerAuc}{0.80}
\newcommand{\PlayerAucCI}{[0.72, 0.88]}
\newcommand{\OppEvents}{274}
\newcommand{\OppGames}{59}
\newcommand{\OppPunished}{0.43}
\newcommand{\OppBrier}{0.145}
\newcommand{\OppBrierBase}{0.247}
\newcommand{\OppBrierSize}{0.231}
\newcommand{\OppBigErrors}{157}
\newcommand{\OppUnpunished}{41}
\newcommand{\OppQPunished}{0.63}
\newcommand{\OppQUnpunished}{0.27}
\newcommand{\OppAuc}{0.81}
\newcommand{\OppAucCI}{[0.73, 0.88]}
\newcommand{\AllPlayerEvents}{334}
\newcommand{\AllPlayerGames}{57}
\newcommand{\AllPlayerPunished}{0.42}
\newcommand{\AllPlayerBrier}{0.162}
\newcommand{\AllPlayerBrierBase}{0.245}
\newcommand{\AllPlayerBrierSize}{0.241}
\newcommand{\AllPlayerBigErrors}{169}
\newcommand{\AllPlayerUnpunished}{43}
\newcommand{\AllPlayerQPunished}{0.57}
\newcommand{\AllPlayerQUnpunished}{0.28}
\newcommand{\AllPlayerAuc}{0.76}
\newcommand{\AllPlayerAucCI}{[0.69, 0.83]}
\newcommand{\AllOppEvents}{355}
\newcommand{\AllOppGames}{60}
\newcommand{\AllOppPunished}{0.43}
\newcommand{\AllOppBrier}{0.151}
\newcommand{\AllOppBrierBase}{0.246}
\newcommand{\AllOppBrierSize}{0.238}
\newcommand{\AllOppBigErrors}{175}
\newcommand{\AllOppUnpunished}{39}
\newcommand{\AllOppQPunished}{0.59}
\newcommand{\AllOppQUnpunished}{0.25}
\newcommand{\AllOppAuc}{0.80}
\newcommand{\AllOppAucCI}{[0.73, 0.86]}
\newcommand{\AppPrior}{0.44}
\newcommand{\PlayerOffset}{250}
\newcommand{\OpponentOffset}{400}
\newcommand{\AllPlayerETwoDiff}{$-$0.001}
\newcommand{\AllPlayerETwoCI}{[$-$0.031, 0.031]}
\newcommand{\CalPlayer}{(0.038,0.086) += (0,0.100) -= (0,0.049) (0.147,0.200) += (0,0.191) -= (0,0.111) (0.312,0.298) += (0,0.142) -= (0,0.111) (0.503,0.567) += (0,0.160) -= (0,0.175) (0.713,0.700) += (0,0.119) -= (0,0.154) (0.909,0.860) += (0,0.074) -= (0,0.133)}
\newcommand{\CalOpp}{(0.042,0.091) += (0,0.105) -= (0,0.051) (0.152,0.184) += (0,0.150) -= (0,0.092) (0.289,0.151) += (0,0.120) -= (0,0.072) (0.513,0.581) += (0,0.135) -= (0,0.148) (0.715,0.757) += (0,0.110) -= (0,0.158) (0.902,0.917) += (0,0.050) -= (0,0.112)}
\newcommand{\CalTable}{%
\begin{tabular}{lrrrrrr}
\toprule
 & \multicolumn{3}{c}{the player's errors} & \multicolumn{3}{c}{the opponents' errors} \\
\cmidrule(lr){2-4}\cmidrule(lr){5-7}
$q$ & mean $q$ & found & $n$ & mean $q$ & found & $n$ \\
\midrule
0.0--0.1 & 0.04 & 0.09 & 58 & 0.04 & 0.09 & 55 \\
0.1--0.2 & 0.15 & 0.20 & 25 & 0.15 & 0.18 & 38 \\
0.2--0.4 & 0.31 & 0.30 & 47 & 0.29 & 0.15 & 53$^{*}$ \\
0.4--0.6 & 0.50 & 0.57 & 30 & 0.51 & 0.58 & 43 \\
0.6--0.8 & 0.71 & 0.70 & 40 & 0.72 & 0.76 & 37 \\
0.8--1.0 & 0.91 & 0.86 & 43 & 0.90 & 0.92 & 48 \\
\bottomrule
\end{tabular}
}
\newcommand{\SoftLo}{0.2}
\newcommand{\SoftHi}{0.4}
\newcommand{\SoftQ}{0.29}
\newcommand{\SoftFound}{0.15}
\newcommand{\SoftN}{53}
\newcommand{\BigLoss}{20}
\newcommand{\BigLossEvents}{58}
\newcommand{\ExampleAQ}{0.954}
\newcommand{\ExampleABand}{1200}
\newcommand{\ExampleALoss}{23}
\newcommand{\ExampleAError}{21\ldots{}\symrook{}cc2}
\newcommand{\ExampleABest}{22.~\symqueen{}xd6}
\newcommand{\ExampleAPlayed}{22.~\symqueen{}xd6}
\newcommand{\ExampleA}{% Written by projects/maia-punish/paper/generate.mjs; do not edit by hand.
\begin{minipage}[t]{0.4\textwidth}\vspace{0pt}
\chessboard[setfen={6k1/3Q1pp1/p2bpqbp/P2p4/3P2P1/5N1P/1rrN1P2/R3R1K1 w - - 2 22}, boardfontsize=13pt, pgfstyle=straightmove, arrow=to, linewidth=0.07em, color=green!50!black, markmoves={d7-d6}]
\end{minipage}\hfill
\begin{minipage}[t]{0.56\textwidth}\vspace{0pt}\small
\textbf{The error:} 21\ldots{}\symrook{}cc2, giving away 22.6\% of the player's winning chances.\par
\textbf{The punishing reply:} 22.~\symqueen{}xd6 (green), the engine's best; Maia-2 at 1200 gives it $q = 0.954$.\par
\textbf{Played:} 22.~\symqueen{}xd6, the punishing reply.\par
\textbf{UCI:} error \uci{c8c2}; punishing reply \uci{d7d6}; played \uci{d7d6}\par
\textbf{FEN:} \fen{6k1/3Q1pp1/p2bpqbp/P2p4/3P2P1/5N1P/1rrN1P2/R3R1K1 w - - 2 22}
\end{minipage}
}
\newcommand{\ExampleBQ}{0.003}
\newcommand{\ExampleBBand}{1300}
\newcommand{\ExampleBLoss}{24}
\newcommand{\ExampleBError}{7.~\symknight{}xd4}
\newcommand{\ExampleBBest}{7\ldots{}e5}
\newcommand{\ExampleBPlayed}{7\ldots{}\symbishop{}b4+}
\newcommand{\ExampleB}{% Written by projects/maia-punish/paper/generate.mjs; do not edit by hand.
\begin{minipage}[t]{0.4\textwidth}\vspace{0pt}
\chessboard[setfen={rnbq1rk1/pp1pbppp/4pn2/8/2PN1B2/3BP3/PP3PPP/RN1QK2R b KQ - 0 7}, boardfontsize=13pt, pgfstyle=straightmove, arrow=to, linewidth=0.07em, color=red!75!black, markmoves={e7-b4}, pgfstyle=straightmove, arrow=to, linewidth=0.07em, color=green!50!black, markmoves={e6-e5}]
\end{minipage}\hfill
\begin{minipage}[t]{0.56\textwidth}\vspace{0pt}\small
\textbf{The error:} 7.~\symknight{}xd4, giving away 24.1\% of the player's winning chances.\par
\textbf{The punishing reply:} 7\ldots{}e5 (green), the engine's best; Maia-2 at 1300 gives it $q = 0.003$.\par
\textbf{Played:} 7\ldots{}\symbishop{}b4+ (red), which gave away 27.1\% of the replier's winning chances.\par
\textbf{UCI:} error \uci{f3d4}; punishing reply \uci{e6e5}; played \uci{e7b4}\par
\textbf{FEN:} \fen{rnbq1rk1/pp1pbppp/4pn2/8/2PN1B2/3BP3/PP3PPP/RN1QK2R b KQ - 0 7}
\end{minipage}
}
\newcommand{\RecomputedIdentical}{matches the records' results byte for byte}
\newcommand{\RerunEvents}{679}
\newcommand{\RecordedEvents}{689}
\newcommand{\AppCommit}{\texttt{3eb8afd}}
\newcommand{\RecordsCommit}{\texttt{248fed3}}
\newcommand{\ReproTable}{%
\begin{tabular}{llrrrrrrll}
\toprule
 & & \multicolumn{2}{c}{events} & \multicolumn{2}{c}{Brier, $q$} & \multicolumn{2}{c}{AUC} & \multicolumn{2}{c}{verdict} \\
\cmidrule(lr){3-4}\cmidrule(lr){5-6}\cmidrule(lr){7-8}\cmidrule(lr){9-10}
replies & to errors of & recorded & rerun & recorded & rerun & recorded & rerun & recorded & rerun \\
\midrule
verified & the player's & 243 & 272 & 0.162 & 0.171 & 0.80 & 0.75 & holds & holds \\
verified & the opponents' & 274 & 307 & 0.145 & 0.158 & 0.81 & 0.83 & holds & holds \\
all & the player's & 334 & 330 & 0.162 & 0.161 & 0.76 & 0.72 & holds & holds \\
all & the opponents' & 355 & 349 & 0.151 & 0.145 & 0.80 & 0.81 & holds & holds \\
\bottomrule
\end{tabular}
}
\newcommand{\RerunBrierMove}{0.014}
\newcommand{\RerunAucMove}{0.05}


\title{Which mistakes get punished at this level?\\[0.3em]
\large A human-move model's odds of the punishing reply, checked on the player's own games}
\author{Chess Coach research lab}
\date{}

\begin{document}
\maketitle

\begin{abstract}
An engine prices a mistake by what it would cost against perfect play. At the player's level,
many mistakes cost nothing, because the reply that punishes them is never found. Maia-2, the human-move model the coaching app ships, gives the probability that a player
at a given rating plays each move. Before building on it, we asked whether its probability of
the punishing reply can be trusted on the player's own games. Over \Events{} errors in \Games{}
games, by a rule fixed before any number was computed, it can: its odds beat both a constant
base rate and a fit on the size of the error, and every bin of its odds is calibrated, for the
player's errors (which the opponents punish) and for the opponents' (which the player punishes),
on verified replies and on all of them. Of the player's errors that cost at least \StepAt{} points
of winning chances, \PlayerUnpunished\% went unpunished, and the odds tell the punished from the
unpunished with an area under the ROC curve of \PlayerAuc. The app now prices every error with
these odds, and since the player-offsets project it reads his opponents at a higher band than
this check did. The lab recomputes the analysis byte for byte from the recorded events, and a
rerun on engine evidence rebuilt on the app at \AppCommit, with each reply priced at the band
that app gives it, reaches the same verdicts.
\end{abstract}

\section{The question}

A review that judges a mistake only by the engine treats every hanging piece alike. At the
player's level that is not how games go: some mistakes are punished at once, and some survive
because the punishing reply is hard to see. A verdict that weighs a mistake \emph{given the
circumstance} needs to know how likely the punishment was. The cheapest circumstance to measure
is the reply itself: did the opponent find the move that punishes the error?

Maia-2 is a neural network trained on human games that gives, for a position and a rating band,
the probability of each legal move. The app runs it in Node through its engine service. If its
probability of the punishing reply is trustworthy on this player's games, a review can say how
often a player at the opponent's level finds the punishment, and the daily coach can rank the
player's mistakes by how likely they are to be punished.

\section{Events}

Every error in the \Games{} games of the accuracy project: a move that gave away at least
\ErrorAt{} points of the mover's winning chances against the engine's best move, read from the
routine-profile reviews. Each error is an event with:

\begin{itemize}[nosep]
\item the \textbf{punishing reply}: the engine's best move in the position after the error;
\item \textbf{E1}: the reply played was exactly that move (the outcome the rule scores);
\item \textbf{E2}: the reply played gave away at most \EAt{} points (an equally good alternative
counts; a secondary outcome);
\item $\boldsymbol{q}$: Maia-2's probability of the punishing reply at the replier's band
(\BandMin{} to \BandMax{} in these games), read from the replier's chess.com rapid rating plus
\PlayerOffset, the app's offset for either player when this check ran.
\end{itemize}

The events run in two directions: the player's errors, which the opponents reply to, and the
opponents' errors, which the player replies to. A reply from the opening book is left out (its
loss is a placeholder). The primary set keeps the \ExactEvents{} replies whose evidence the
review's deeper pass verified, which is how the question reads ``the engine's best reply''.
Because that pass checks errors first, the set is selected partly on the outcome, so every result
is also given on all \Events{} non-book replies.

\section{The rule}

Fixed before any number was computed: Maia is trusted in a direction when (1) its Brier score for
E1 beats both a constant base rate and a logistic regression on the error's size alone, both
fitted leaving one game out, with the 95\% interval of each difference below zero over \Draws{}
bootstrap draws of games; and (2) no calibration bin of $q$ with \BinMin{} or more events has its
mean $q$ more than \BinGap{} outside the Wilson interval of the share that was found.

\section{Results}

The rule holds in both directions and on both sets (Table~\ref{tab:results}). On the player's
errors, Maia's Brier score is \PlayerBrier{} against \PlayerBrierBase{} for the base rate and
\PlayerBrierSize{} for the fit on the error's size; on the opponents' errors, \OppBrier{} against
\OppBrierBase{} and \OppBrierSize. The secondary outcome is weaker: on all replies to the player's
errors, Maia does not beat the base rate for E2 (difference \AllPlayerETwoDiff,
\AllPlayerETwoCI), since an equally good alternative is often a different move.

\begin{table}[!htbp]
\centering\scriptsize
\setlength{\tabcolsep}{3.5pt}
\ResultTable
\caption{Brier scores for E1 (the punishing reply is played): Maia's odds ($q$), the base rate
(\emph{base}) and the fit on the error's size (\emph{size}), both leaving one game out, and the
paired differences with 95\% bootstrap intervals over games. Events with games in parentheses;
\emph{punished} is the share of E1.}
\label{tab:results}
\end{table}

\paragraph{Calibration.} Every bin is inside the rule (Table~\ref{tab:cal},
Figure~\ref{fig:cal}). One bin is flagged by a stricter reading (its mean $q$ outside the Wilson
interval and more than \BinGap{} from what was found): the player punishes the opponents' errors
less often than Maia expects where $q$ is \SoftLo{} to \SoftHi{} (\SoftFound{} found against a
mean $q$ of \SoftQ, $n = \SoftN$). It is inside the rule, and worth watching as games accumulate.

\begin{table}[!htbp]
\centering\small
\CalTable
\caption{Calibration on verified replies: by bin of $q$, its mean, the share of punishing replies
found, and the number of events. $^{*}$ flagged by the stricter reading.}
\label{tab:cal}
\end{table}

\begin{figure}[!htbp]
\centering
\begin{tikzpicture}
\begin{axis}[width=0.62\textwidth, height=0.62\textwidth, xmin=0, xmax=1, ymin=0, ymax=1,
  xlabel={Maia-2's odds of the punishing reply, mean $q$ in the bin},
  ylabel={share of punishing replies found}, grid=major, legend pos=north west,
  legend cell align=left]
\addplot[gray, densely dotted, domain=0:1] {x};
\addlegendentry{calibrated}
\addplot+[only marks, mark=*, mark size=1.6pt, black, error bars/.cd, y dir=both, y explicit]
  coordinates {\CalPlayer};
\addlegendentry{the player's errors}
\addplot+[only marks, mark=square*, mark size=1.6pt, red!70!black, error bars/.cd, y dir=both,
  y explicit] coordinates {\CalOpp};
\addlegendentry{the opponents' errors}
\end{axis}
\end{tikzpicture}
\caption{Calibration on verified replies, with 95\% Wilson intervals.}
\label{fig:cal}
\end{figure}

\section{What it means for verdicts}

Of the errors that cost at least \StepAt{} points, \PlayerUnpunished\% of the player's
(\PlayerBigErrors{} errors) and \OppUnpunished\% of the opponents' (\OppBigErrors) went
unpunished in these games: the reply played gave away more than \EAt{} points. The odds tell
them apart. The mean $q$ was \PlayerQPunished{} for the player's errors that were punished and
\PlayerQUnpunished{} for those that were not, and the area under the ROC curve is \PlayerAuc{}
\PlayerAucCI{} for the player's errors and \OppAuc{} \OppAucCI{} for the opponents' (on all
replies, \AllPlayerAuc{} and \AllOppAuc).

Figures~\ref{fig:exa} and~\ref{fig:exb} show two of the player's errors that cost at least
\BigLoss{} points, on verified replies (\BigLossEvents{} such errors): the punished one Maia was surest
of, other than taking back the piece that just moved, and the unpunished one Maia was surest
would be missed. The green arrow is the punishing reply; the red one, the reply played when it
was something else.

\begin{figure}[!htbp]\ExampleA\caption{Found: Maia-2 gives \ExampleABest{} $q = \ExampleAQ$
at \ExampleABand, and the opponent played it.}\label{fig:exa}\end{figure}
\begin{figure}[!htbp]\ExampleB\caption{Missed: Maia-2 gives \ExampleBBest{} $q = \ExampleBQ$
at \ExampleBBand, and the opponent played \ExampleBPlayed{} instead.}\label{fig:exb}\end{figure}

\paragraph{In the app.} The review service now prices every error of at least \ErrorAt{} points
with these odds, computed as here (the replier's band, the engine's best answer), with one change
since: the player-offsets project found that Maia-2 predicts his opponents' moves better with
\OpponentOffset{} added to their ratings, so the app reads an opponent's reply at that band and
keeps \PlayerOffset{} for the player's own replies. Game Review
shows the line under a mistake, and the routine's teaching pass picks its deep looks by the
winning chances lost times the odds. A turning point with no odds of its own is ranked at
\AppPrior, the share of the player's errors punished in this check.

\section{Reproduction in the lab}

\paragraph{Recomputed.} The lab runs the analysis (\texttt{analyze.mjs}) over the events the
records keep, and its results \RecomputedIdentical. The records are read at \RecordsCommit.

\paragraph{Rerun.} The lab rebuilt the routine evidence for the same games on the app at
\AppCommit, found the errors and their punishing replies again, and ran Maia-2 on each at the
band that app gives the replier: an opponent's rating plus \OpponentOffset{} after the player's
error, the player's plus \PlayerOffset{} after an opponent's. Every band and probability matches
the odds the app's own review recorded for that reply. The analysis then ran on \RerunEvents{}
events against \RecordedEvents{} recorded, since the rebuilt evidence differs in detail. Every
verdict holds again; no Brier score moves by more than
\RerunBrierMove{} and no area under the curve by more than \RerunAucMove{}
(Table~\ref{tab:repro}).

\begin{table}[!htbp]
\centering\small
\setlength{\tabcolsep}{4pt}
\ReproTable
\caption{The recorded analysis beside the rerun on rebuilt evidence. AUC is for E2 on errors of
at least \StepAt{} points.}
\label{tab:repro}
\end{table}

\section{Limitations}

\begin{itemize}[nosep]
\item \Games{} games of one player and their opponents, all near the same rating. The check says
Maia can be trusted at this level, not at every level.
\item E1 counts only the engine's best move as the punishment. A reply that punishes as well by
another route counts as missed under E1 and found under E2.
\item The band comes from a conversion of the chess.com rating, and a player's form on the day is
not in it.
\end{itemize}

\end{document}
